\section{Ethical Considerations}
\label{sec:ethics}

Publication venues increasingly expect an ethics statement; \citet{benotti-blackburn-2022-ethics} classify such sections by whether they address research benefits, potential harms and affected vulnerable groups, and \citet{karamolegkou-etal-2025-ethical} catalogue the concerns raised across ACL papers.
This section addresses all three, for the original work and for the new data.

\textbf{Responsible use of uncertainty.}
The original NGBoost work makes predictive uncertainty available alongside point estimates \cite{duan2020ngboost}, potentially supporting better-informed decisions.
However, a predictive distribution can create misplaced confidence when its assumptions fail or its intervals are interpreted as guarantees.
The Superconductivity results illustrate why coverage alone is insufficient: high coverage accompanied much wider intervals and worse point predictions (Section~\ref{sec:new-results}).
In health or financial applications, misleading uncertainty could affect treatment or valuation decisions.
Neither the original benchmark results nor this extension establishes suitability for such deployment; calibration and usefulness require assessment in the intended setting.

\textbf{Benchmark data in the original work.}
The original regression benchmark includes the Boston housing dataset (Section~\ref{sec:original-data}).
That dataset carries a constructed variable derived from the proportion of Black residents in each town, built on an assumption that racial segregation raised house prices; the scikit-learn maintainers removed its loader for this reason \cite{sklearnboston}.
The original paper reports an averaged benchmark result over that dataset without discussing how it was built.
A paper published on this data today would be expected to state such provenance; the Boston data are not used in this project.

\textbf{Licensing, ownership and attribution.}
Open access supports replication without transferring ownership of reused material.
NGBoost's Apache 2.0 licence requires distributed copies to retain relevant notices and the licence, and modified upstream files to identify changes \cite{ngboostrepo}.
UCI lists the five additional public datasets (Section~\ref{sec:new-data}) under CC BY 4.0, which requires credit, a licence link and an indication of changes when the data or adaptations are redistributed \cite{creativecommons4}.
Section~\ref{sec:new-data} identifies their creators and preparation steps.
Yacht is likewise listed under CC BY 4.0, while no separate licence statement was located for Kin8nm, which reached this project through the NGBoost repository; it is therefore cited rather than redistributed.
The judge uses a third-party quantisation of the Apache 2.0 MiniCPM4.1-8B release \cite{openbmb2025minicpm41}, referenced by repository and checksum rather than redistributed (Appendix~\ref{app:new-results-details}).
The public repository released with this report (Section~\ref{sec:paper}) separates upstream material from project-authored files and keeps third-party terms on the upstream files rather than placing everything under one project licence.
These practices follow three principles of the Australian Code for the Responsible Conduct of Research: fairness, which includes citing and crediting the work of others; transparency in reporting methods, data and findings; and accountability \cite{australiancode2018}.

\textbf{Data sensitivity and generalisation.}
Parkinsons contains health-related measurements, while Real Estate includes transaction locations.
Permission to reuse data does not remove privacy considerations or establish consent for every downstream application.
Parkinsons subject identifiers define disjoint partitions rather than predictors; this reduces evaluation leakage but does not itself establish privacy protection.
Aggregate results do not demonstrate equal performance across demographic groups, populations or locations, and subgroup fairness was not evaluated.
The synthetic experiments use no personal records, and because their generating distributions are known, the predicted uncertainty can be checked against the truth.
Those outcomes are Normal given the inputs, with a specified mean and scale, so the synthetic results do not establish how the method would behave on real clinical or property data.

\textbf{Transparent reporting and automated judgement.}
Preserved outputs, fixed protocols, seeds and checksums make the reported comparisons open to scrutiny.
Reporting the published results, the released-code run and the controlled selector experiment as separate evidence lines avoids implying stronger replication evidence than was obtained.
The explanation judge rejected all five texts that misdescribe the prediction interval when scoring them individually, yet preferred those texts whenever they appeared second in a side-by-side comparison (Section~\ref{sec:new-results}).
Its agreement scores therefore do not justify autonomous explanation selection, and anyone using it should keep human review of statistical interpretation and evidential support.
